\section{Methodology}

The workflow begins with a simplified finite-gain, finite-bandwidth behavioural TIA model. MATLAB scripts generate a baseline response and then sweep selected parameters. The baseline run compares ideal and finite-gain responses using the repository baseline script, with outputs stored in the baseline response and metrics CSV files (Figure~\ref{fig:baseline-response}; Table~\ref{tab:model-assumptions}).

\ModelAssumptionsTable

\begin{figure}[H]
  \centering
  \begin{subfigure}[t]{0.48\linewidth}
    \centering
    \includegraphics[width=\linewidth]{figures/tia_baseline_magnitude.png}
    \caption{Magnitude response.}
  \end{subfigure}
  \hfill
  \begin{subfigure}[t]{0.48\linewidth}
    \centering
    \includegraphics[width=\linewidth]{figures/tia_baseline_phase.png}
    \caption{Phase response.}
  \end{subfigure}
  \caption{Baseline MATLAB behavioural transimpedance response for the selected photodiode TIA case. Panel (a) shows the magnitude response and panel (b) shows the inversion-removed phase response under ideal and finite-gain op-amp assumptions. The curves are behavioural model outputs and are not vendor macromodel or hardware measurement data.}
  \label{fig:baseline-response}
\end{figure}
\FloatBarrier

The feedback-capacitance sweep exposes the bandwidth and peaking trade-off. Within the selected behavioural sweep, larger $C_f$ reduces extracted bandwidth and generally lowers peaking. This result is interpreted within the explored parameter range, not as a universal compensation law (Figures~\ref{fig:cf-bandwidth} and~\ref{fig:cf-peaking}).

\begin{figure}[H]
  \centering
  \includegraphics[width=0.75\linewidth]{figures/tia_bandwidth_vs_Cf.png}
  \caption{Extracted $-3$ dB bandwidth versus feedback capacitance for the selected MATLAB behavioural TIA sweep. The figure shows the bandwidth cost associated with increasing $C_f$ under the explored assumptions; it is not a universal compensation rule.}
  \label{fig:cf-bandwidth}
\end{figure}

\begin{figure}[H]
  \centering
  \includegraphics[width=0.75\linewidth]{figures/tia_peaking_vs_Cf.png}
  \caption{Extracted gain peaking versus feedback capacitance for the selected MATLAB behavioural TIA sweep. Peaking is used here as a project screening metric and stability proxy, not as a complete proof of phase margin or hardware stability.}
  \label{fig:cf-peaking}
\end{figure}
\FloatBarrier

The design-region stage maps selected $R_f$, $C_{pd}$, $C_f$, and $f_t$ values into project-defined Safe/Marginal/Risky categories. The criteria and labels are summarized explicitly as project-defined screening definitions (Table~\ref{tab:screening-criteria}). The region map and representative response examples provide the main visual evidence for this stage (Figures~\ref{fig:design-region-map} and~\ref{fig:representative-responses}).

\ScreeningCriteriaTable

\begin{figure}[H]
  \centering
  \includegraphics[width=0.78\linewidth]{figures/tia_design_region_map_Cpd_ft.png}
  \caption{Project-defined behavioural design-region map over photodiode capacitance and op-amp transition frequency. Each tile selects the best available feedback capacitance from the behavioural sweep according to the repository's Safe/Marginal/Risky screening labels. The labels organize the explored design space and are not universal stability classifications.}
  \label{fig:design-region-map}
\end{figure}

\begin{figure}[H]
  \centering
  \includegraphics[width=0.78\linewidth]{figures/tia_representative_region_responses.png}
  \caption{Representative MATLAB behavioural transimpedance responses for selected project-defined Safe, Marginal, and Risky cases. The figure connects the screening labels to example frequency-response shapes within the explored assumptions; it does not certify hardware stability.}
  \label{fig:representative-responses}
\end{figure}
\FloatBarrier

\subsection{Behavioural TIA Model}

The behavioural model represents the photodiode TIA as a feedback transimpedance stage with feedback resistance $R_f$, feedback capacitance $C_f$, photodiode capacitance $C_{pd}$, low-frequency op-amp gain $A_0$, and op-amp transition frequency $f_t$. The implementation computes the response, and the metric extraction stage extracts low-frequency gain, bandwidth, peaking, and phase-related metrics. Table~\ref{tab:model-assumptions} lists the baseline assumptions and swept variables used to produce the manuscript figures.

This model makes parameter dependence explicit enough for rapid sweep generation. It shows how selected $C_f$, $C_{pd}$, and $f_t$ assumptions influence bandwidth and peaking, and it identifies regions that deserve later vendor macromodel or circuit-level review. The model remains a first-order pre-design approximation. It does not reproduce detailed vendor macromodel dynamics, transistor-level behaviour, layout-aware extraction, or measurement. Full-text TIA literature motivates the variables and trade-offs, but the project does not reproduce the circuit methods of Analui 2004, Lu 2007, Noh 2020, or regulated-cascode CMOS TIA papers \citep{ANALUI_2004_BW_ENHANCEMENT,LU_2007_BROADBAND_TIA,NOH_2020_CF_TIA_DC_LOOP,PARK_YOO_2004_RGC_TIA}.

Table~\ref{tab:model-assumptions} is kept as one combined manuscript table in this draft. It combines baseline assumptions, focused feedback-capacitance sweep ranges, full behavioural sweep counts, and design-region evidence layers, and the final version will identify which evidence layer supports each claim after target-journal formatting.

\subsection{Design-Region Mapping}

The design-region workflow translates behavioural sweep metrics into project-defined categories. Safe, Marginal, and Risky are internal screening labels, not universal stability classifications and not hardware certification labels. They organize the current sweep results and help select cases for review. Peaking is used as a practical screening metric in this workflow, but it is not by itself a complete phase-margin or stability proof. The map is therefore useful for prioritizing further analysis, not for certifying hardware behaviour.

\subsection{Vendor Macromodel Comparison}

The vendor macromodel comparison stage adds model-specific simulation context. It compares selected behavioural trends with OP27, OPA818, and ADA4817 vendor macromodel export data under the repository's stated test assumptions. These comparisons are not hardware measurements and do not constitute experimental validation.

The OP27 smoke-test set contains three feedback-capacitance cases. In the repository evidence, the smallest-$C_f$ OP27 case exhibits strong peaking, while larger-$C_f$ cases reduce peaking and bandwidth. This result supports a qualitative compensation caution in a low-GBW macromodel case. OP27 is treated as a smoke-test comparison point rather than a preferred photodiode TIA recommendation.

The OPA818 set adds four high-speed vendor macromodel cases. In the selected OPA818 vendor macromodel rows, the tested cases remain project-defined Safe under the peaking metric and stated test assumptions. The result provides selected macromodel evidence under the specified supply, capacitance, and export assumptions, not a proof that OPA818 is universally stable or optimal.

The ADA4817 set adds a third vendor macromodel comparison. In the selected ADA4817 vendor macromodel rows, the smallest-$C_f$ case is project-defined Marginal under the peaking metric and stated test assumptions, while the larger-$C_f$ cases are project-defined Safe under the same table context. This contrast supports the broader observation that vendor macromodels can show different bandwidth and peaking boundaries under comparable project test structures.

\VendorSummaryTable

\begin{figure}[H]
  \centering
  \includegraphics[width=0.78\linewidth]{figures/spice_vendor_model_bandwidth_peaking_summary.png}
  \caption{Three-vendor macromodel comparison summary for OP27, OPA818, and ADA4817 feedback-capacitance cases. The plot compares extracted peaking and $-3$ dB bandwidth from vendor macromodel exports under the repository's test assumptions. This is vendor macromodel comparison evidence, not hardware validation or complete SPICE coverage.}
  \label{fig:vendor-summary}
\end{figure}
\FloatBarrier

\subsection{First-Pass Noise Analysis}

The first-pass noise workflow extends the bandwidth and peaking discussion by estimating relative noise contributions under configured assumptions. The model includes configured feedback resistor, op-amp voltage-noise, op-amp current-noise, and photodiode shot-noise terms. The baseline contribution figure summarizes the configured output-noise terms for the baseline behavioural case (Figure~\ref{fig:noise-baseline}).

These outputs are first-pass behavioural estimates. They are not measured noise, detector-complete noise validation, or SPICE noise analysis. Their purpose is to show how the same workflow can carry noise assumptions into the screening discussion when compensation choices change. The noise-bandwidth trade-off figure connects this estimate back to the compensation sweep (Figure~\ref{fig:noise-tradeoff}).

The literature supports the relevance of discussing noise and bandwidth together. Noh 2020 supports capacitive-feedback and noise/stability context \citep{NOH_2020_CF_TIA_DC_LOOP}. Lu 2007 and Abdollahi 2025 support broad-band and noise-aware TIA framing \citep{LU_2007_BROADBAND_TIA,ABDOLLAHI_2025_MDFRGC_TIA}.
% TODO citation pending manual check: LI_2021_LOW_NOISE_OPTICAL_RX_TIA is excluded from references_draft.bib.
